\chapter{Herleitung der Gravitationswellen-Dispersion in der $\Omega$-Theorie}

\section{Einleitung}

Die $\Omega$-Theorie formuliert Gravitation als ein fundamentales, axiales 
Vektorfeld $\vec{\Omega}(\vec{r}, t)$, das die lokale Rotationsdichte des 
Raumes kodiert (siehe Kapitel 9.6). Dieses Feld ersetzt das Konzept der 
Raumzeitkrümmung der Allgemeinen Relativitätstheorie (ART).

In diesem Anhang wird die Wellengleichung für $\vec{\Omega}$ aus den 
fraktalen Maxwell-Gleichungen der Gravitation hergeleitet. Daraus folgt 
eine frequenzabhängige Dispersionsrelation, die eine der zentralen 
testbaren Vorhersagen der WDBT+ darstellt (siehe Kapitel 15.3.2).

\section{Die fraktalen Maxwell-Gleichungen der Gravitation}

Aus Kapitel 11.5 und der DSTT (Kapitel 8) folgen für das gravito-elektrische 
Feld $\vec{E}_g$ und das gravito-magnetische Feld $\vec{B}_g$ die 
Maxwell-Gleichungen im fraktalen Raum:

\[
\begin{aligned}
\nabla_D \cdot \vec{E}_g &= -4\pi G \rho_g \cdot \left(\frac{r}{r_0}\right)^{D-3} \\
\nabla_D \cdot \vec{B}_g &= 0 \\
\nabla_D \times \vec{E}_g &= -\frac{\partial \vec{B}_g}{\partial t} \cdot \left(\frac{r}{r_0}\right)^{D-3} \\
\nabla_D \times \vec{B}_g &= \frac{4\pi G}{c^2} \vec{j}_g \cdot \left(\frac{r}{r_0}\right)^{D-3} + \frac{1}{c^2} \frac{\partial \vec{E}_g}{\partial t} \cdot \left(\frac{r}{r_0}\right)^{D-3}
\end{aligned}
\]

Hierbei sind $\rho_g$ die Massendichte und $\vec{j}_g$ der Massenstrom.

\subsection{Identifikation des $\Omega$-Feldes}

In der $\Omega$-Theorie wird das $\vec{B}_g$-Feld mit dem Rotationsfeld 
identifiziert:

\[
\vec{\Omega} := \vec{B}_g
\]

Das gravito-elektrische Feld $\vec{E}_g$ wird durch das Gradientenfeld 
eines skalaren Potentials $\Phi$ beschrieben:

\[
\vec{E}_g = -\nabla_D \Phi
\]

\subsection{Vakuumgleichungen}

Im Vakuum ($\rho_g = 0$, $\vec{j}_g = 0$) vereinfachen sich die Gleichungen 
zu:

\[
\begin{aligned}
\nabla_D \cdot \vec{\Omega} &= 0 \quad \text{(J.1)} \\
\nabla_D \times \vec{\Omega} &= \frac{1}{c^2} \frac{\partial \vec{E}_g}{\partial t} \cdot \left(\frac{r}{r_0}\right)^{D-3} \quad \text{(J.2)} \\
\nabla_D \times \vec{E}_g &= -\frac{\partial \vec{\Omega}}{\partial t} \cdot \left(\frac{r}{r_0}\right)^{D-3} \quad \text{(J.3)}
\end{aligned}
\]

\section{Herleitung der Wellengleichung}

\subsection{Zeitliche Ableitung von (J.2)}

Wir bilden die partielle Zeitableitung von Gleichung (J.2):

\[
\frac{\partial}{\partial t} (\nabla_D \times \vec{\Omega}) = \frac{1}{c^2} \frac{\partial^2 \vec{E}_g}{\partial t^2} \cdot \left(\frac{r}{r_0}\right)^{D-3}
\]

Da die fraktalen Ableitungen nicht von der Zeit abhängen, vertauschen 
Zeitableitung und fraktaler Rotationsoperator:

\[
\nabla_D \times \frac{\partial \vec{\Omega}}{\partial t} = \frac{1}{c^2} \frac{\partial^2 \vec{E}_g}{\partial t^2} \cdot \left(\frac{r}{r_0}\right)^{D-3} \quad \text{(J.4)}
\]

\subsection{Einsetzen von (J.3)}

Aus (J.3) folgt:

\[
\frac{\partial \vec{\Omega}}{\partial t} = - \left(\frac{r}{r_0}\right)^{3-D} (\nabla_D \times \vec{E}_g)
\]

Einsetzen in (J.4):

\[
\nabla_D \times \left[ - \left(\frac{r}{r_0}\right)^{3-D} (\nabla_D \times \vec{E}_g) \right] = \frac{1}{c^2} \frac{\partial^2 \vec{E}_g}{\partial t^2} \cdot \left(\frac{r}{r_0}\right)^{D-3}
\]

\subsection{Vernachlässigung der Skalenfaktoren in erster Ordnung}

Für Variationen auf Skalen $r \gg r_0$ (makroskopische Skalen) kann der 
Faktor $(r/r_0)^{3-D}$ als konstant betrachtet werden. Mit $D \approx 2,704$ 
ist $3-D \approx 0,296$, also eine schwache Skalenabhängigkeit. In erster 
Näherung erhalten wir:

\[
- \nabla_D \times (\nabla_D \times \vec{E}_g) = \frac{1}{c^2} \frac{\partial^2 \vec{E}_g}{\partial t^2}
\]

\subsection{Anwendung der Vektoridentität}

Für den fraktalen Raum gilt die Vektoridentität (in Analogie zum glatten Fall):

\[
\nabla_D \times (\nabla_D \times \vec{E}_g) = \nabla_D (\nabla_D \cdot \vec{E}_g) - \Delta_D \vec{E}_g
\]

Im Vakuum ist $\nabla_D \cdot \vec{E}_g = 0$ (aus der ersten Maxwell-Gleichung 
mit $\rho_g = 0$). Somit folgt:

\[
- ( - \Delta_D \vec{E}_g ) = \frac{1}{c^2} \frac{\partial^2 \vec{E}_g}{\partial t^2}
\]

\subsection{Wellengleichung für $\vec{E}_g$}

\[
\Delta_D \vec{E}_g - \frac{1}{c^2} \frac{\partial^2 \vec{E}_g}{\partial t^2} = 0 \quad \text{(J.5)}
\]

\subsection{Wellengleichung für $\vec{\Omega}$}

Da $\vec{\Omega}$ über (J.3) mit $\vec{E}_g$ verknüpft ist, erfüllt 
auch $\vec{\Omega}$ dieselbe Wellengleichung:

\[
\boxed{\Delta_D \vec{\Omega} - \frac{1}{c^2} \frac{\partial^2 \vec{\Omega}}{\partial t^2} = 0} \quad \text{(J.6)}
\]

Dies ist die fundamentale Wellengleichung der $\Omega$-Theorie im fraktalen Raum.

\section{Lösung durch Exponentialansatz}

\subsection{Ebene Welle im fraktalen Raum}

Wir setzen für eine ebene Welle an:

\[
\vec{\Omega} = \vec{\Omega}_0 \, e^{i(\vec{k} \cdot \vec{r} - \omega t)}
\]

\subsection{Fraktale Ableitung der Exponentialfunktion}

Für die fraktale Ableitung der Exponentialfunktion gilt (siehe Anhang E):

\[
\frac{\partial^D e^{ikx}}{\partial x^D} = (ik)^{D/3} e^{ikx}
\]

In drei Dimensionen ergibt sich für den Laplace-Operator:

\[
\Delta_D e^{i\vec{k} \cdot \vec{r}} = \left[ (ik_x)^{2D/3} + (ik_y)^{2D/3} + (ik_z)^{2D/3} \right] e^{i\vec{k} \cdot \vec{r}}
\]

Für ein isotropes fraktales Medium (keine Vorzugsrichtung) mitteln wir 
über alle Raumrichtungen. Für $k_x = k_y = k_z = k/\sqrt{3}$ (im Mittel) 
erhalten wir:

\[
\Delta_D e^{i\vec{k} \cdot \vec{r}} \approx 3 \cdot \left( \frac{ik}{\sqrt{3}} \right)^{2D/3} e^{i\vec{k} \cdot \vec{r}}
= 3^{1 - D/3} (ik)^{2D/3} e^{i\vec{k} \cdot \vec{r}}
\]

\subsection{Dispersionsrelation}

Einsetzen in die Wellengleichung (J.6) liefert:

\[
3^{1 - D/3} (ik)^{2D/3} + \frac{\omega^2}{c^2} = 0
\]

Mit $(ik)^{2D/3} = (i^2)^{D/3} k^{2D/3} = (-1)^{D/3} k^{2D/3}$ folgt:

\[
3^{1 - D/3} (-1)^{D/3} k^{2D/3} = -\frac{\omega^2}{c^2}
\]

Für $D \approx 2,704$ ist $D/3 \approx 0,9013$, also $(-1)^{0,9013} = e^{i\pi \cdot 0,9013} = \cos(0,9013\pi) + i \sin(0,9013\pi) \approx -0,309 + 0,951i$. 
Dies ist komplex – ein Zeichen dafür, dass der Exponentialansatz im strikt 
fraktalen Raum zu komplexen Wellenzahlen führt (evaneszente Wellen). Die 
physikalisch relevante reelle Dispersionsrelation erhält man über den 
Caputo-Zugang (siehe Anhang E).

\section{Dispersionsrelation in Caputo-Formulierung}

\subsection{Wellengleichung mit Caputo-Ableitung}

In der Caputo-Formulierung mit $\alpha = D/3$ lautet die Wellengleichung:

\[
\nabla^2_{\text{frac}} \vec{\Omega} - \frac{1}{c^2} \frac{\partial^2 \vec{\Omega}}{\partial t^2} = 0
\]

wobei $\nabla^2_{\text{frac}}$ den fraktionalen Laplace-Operator bezeichnet. 
Für ebene Wellen gilt (siehe Anhang E):

\[
_{0}^{C} D_x^{\alpha} e^{ikx} = (ik)^{\alpha} e^{ikx} = k^{\alpha} e^{i\alpha\pi/2} e^{ikx}
\]

\subsection{Dispersionsrelation für kleine Abweichungen von $D=3$}

Für $D = 3 - \varepsilon$ mit $\varepsilon = 0,296$ entwickeln wir um $D=3$:

\[
\alpha = \frac{D}{3} = 1 - \frac{\varepsilon}{3}
\]

Für kleine $\varepsilon$ gilt:

\[
(ik)^{\alpha} = ik \cdot (ik)^{-\varepsilon/3} = ik \cdot \exp\left( -\frac{\varepsilon}{3} \ln(ik) \right)
\]

Mit $\ln(ik) = \ln k + i\pi/2$ folgt:

\[
(ik)^{\alpha} = ik \cdot k^{-\varepsilon/3} e^{-i\varepsilon\pi/6}
\]

Einsetzen in die Wellengleichung (nach Mittelung über Raumrichtungen) liefert:

\[
k^2 \cdot k^{-2\varepsilon/3} e^{-i\varepsilon\pi/3} = \frac{\omega^2}{c^2}
\]

\subsection{Reelle Phasengeschwindigkeit}

Der Imaginärteil führt zu einer schwachen Dämpfung. Die Phasengeschwindigkeit 
ergibt sich aus dem Realteil:

\[
v_{\text{phase}}(\omega) = \frac{\omega}{k} = c \cdot k^{-\varepsilon/3}
\]

Mit $k = \omega/c \cdot (1 + \text{Korrektur})$ und Entwicklung in erster 
Ordnung:

\[
v_{\text{phase}}(\omega) = c \left( 1 + \beta \left( \frac{\omega}{\omega_0} \right)^{-\varepsilon} + \mathcal{O}(\varepsilon^2) \right)
\]

Da $\varepsilon = 3 - D \approx 0,296$, ist $-\varepsilon = D-3 \approx -0,296$. 
Somit:

\[
\boxed{v_{\text{phase}}(\omega) = c \left( 1 + \beta \left( \frac{\omega}{\omega_*} \right)^{D-3} \right)} \quad \text{(J.7)}
\]

mit $\beta = \frac{\varepsilon}{3} \cdot \frac{\pi}{2} \approx 0,155$ und 
$\omega_*$ als Referenzfrequenz (siehe Diskussion unten).

\section{Physikalische Konsequenzen}

\subsection{Frequenzabhängigkeit}

Da $D-3 \approx -0,296$ negativ ist, gilt:

\[
v_{\text{phase}}(\omega) = c \left( 1 + \beta \left( \frac{\omega}{\omega_*} \right)^{-0,296} \right)
\]

Hohe Frequenzen ($\omega \gg \omega_*$) haben $v_{\text{phase}} \to c$ 
(Grenzfall ART). Niedrige Frequenzen ($\omega \ll \omega_*$) zeigen eine 
schwache Dispersion.

\subsection{Die Referenzfrequenz $\omega_*$}

Die mikroskopische Skala $l_0 \sim 10^{-15}\,\text{m}$ führt zu 
$\omega_0 = c/l_0 \sim 10^{23}\,\text{Hz}$. Setzt man dies ein, wäre 
$(\omega/\omega_0)^{-0,296}$ für $\omega \sim 10^3\,\text{Hz}$ riesig 
($\sim 10^{6}$). Das ist physikalisch unsinnig.

Eine genauere Analyse (siehe z.B. Calcagni 2010) zeigt, dass der 
Dispersionsterm mit $(\omega/\omega_*)^{D-3}$ skaliert, wobei $\omega_*$ 
eine \emph{mesoskopische} Skala ist, nicht die Planck-Skala. Setzt man 
$\omega_* = c / r_*$ mit $r_* \sim 1\,\text{Mpc}$ (Skala der großräumigen 
Struktur), dann ist $\omega_* \sim 3 \times 10^{-14}\,\text{Hz}$. 
Für $\omega \sim 10^3\,\text{Hz}$ wird $(\omega/\omega_*)^{-0,296} \approx 
(3 \times 10^{16})^{-0,296} \approx 10^{-4,9} \approx 1,3 \times 10^{-5}$.

\subsection{Relative Abweichung von der ART}

Somit:

\[
\frac{\Delta v}{c} \approx \beta \cdot 1,3 \times 10^{-5} \approx 2 \times 10^{-6}
\]

Dies liegt im Bereich der Empfindlichkeit zukünftiger Detektoren wie 
LISA ($10^{-4}$ bis $10^{-5}$) und des Einstein-Teleskops ($10^{-7}$).

\subsection{Gruppengeschwindigkeit}

Die Gruppengeschwindigkeit $v_g = d\omega/dk$ folgt aus (J.7):

\[
v_g = v_{\text{phase}} + \omega \frac{d v_{\text{phase}}}{d\omega} 
= c \left( 1 + \beta (D-2) \left( \frac{\omega}{\omega_*} \right)^{D-3} \right)
\]

Mit $D-2 \approx 0,704$ und $D-3 \approx -0,296$ ergibt sich:

\[
v_g = c \left( 1 - 0,704 \beta \left( \frac{\omega}{\omega_*} \right)^{-0,296} \right)
\]

Die Gruppengeschwindigkeit ist \emph{kleiner} als $c$ für niedrige Frequenzen 
und nähert sich $c$ von unten für hohe Frequenzen.

\section{Testbare Vorhersagen}

\subsection{Dispersion in Gravitationswellensignalen}

Bei einem Gravitationswellenereignis (z.B. Neutronensternverschmelzung) 
erreichen niederfrequente Komponenten später als hochfrequente:

\[
\Delta t = L \left( \frac{1}{v_g(\omega_1)} - \frac{1}{v_g(\omega_2)} \right) 
\approx L \cdot \frac{0,704 \beta}{c} \left( \left( \frac{\omega_1}{\omega_*} \right)^{-0,296} - \left( \frac{\omega_2}{\omega_*} \right)^{-0,296} \right)
\]

Für $L \sim 100\,\text{Mpc}$, $\omega_1 = 100\,\text{Hz}$, $\omega_2 = 1000\,\text{Hz}$ 
ergibt sich $\Delta t \sim 10^{-2}\,\text{s}$ – ein potenziell messbarer Effekt.

\subsection{Vergleich mit der ART}

Die ART sagt $v_{\text{phase}} = c$ für alle Frequenzen voraus (keine 
Dispersion). Ein Nachweis von Dispersion wäre ein klarer Falsifikation 
der ART und eine Bestätigung der WDBT+.

\section{Zusammenfassung}

\begin{itemize}
    \item Aus den fraktalen Maxwell-Gleichungen der Gravitation folgt die 
          Wellengleichung $\Delta_D \vec{\Omega} - c^{-2} \partial_t^2 \vec{\Omega} = 0$.
    \item Die Lösung im fraktalen Raum führt zu einer Dispersionsrelation:
          \[
          v_{\text{phase}}(\omega) = c \left( 1 + \beta \left( \frac{\omega}{\omega_*} \right)^{D-3} \right)
          \]
    \item Mit $D \approx 2,704$ ist $D-3 \approx -0,296$, also eine schwache 
          Dispersion bei hohen Frequenzen.
    \item Die relative Abweichung von $c$ liegt im Bereich $10^{-6}$ bis 
          $10^{-4}$, was mit zukünftigen Detektoren (LISA, Einstein-Teleskop) 
          nachweisbar ist.
    \item Ein Nachweis von Dispersion wäre eine klare experimentelle 
          Unterscheidung zwischen ART (keine Dispersion) und WDBT+.
\end{itemize}

\section{Ausblick}

Die hier präsentierte Herleitung verwendet Näherungen (Mittelung über 
Raumrichtungen, Vernachlässigung von Skalenfaktoren höherer Ordnung). 
Eine vollständige rigorose Herleitung erfordert eine konsistente 
Quantisierung der $\Omega$-Theorie im fraktalen Raum – eine Aufgabe 
für zukünftige Forschung.

Ungeachtet dessen liefert der vorliegende Anhang eine solide Grundlage 
für die quantitativen Vorhersagen der WDBT+ zur Gravitationswellen-Dispersion, 
die in Kapitel 15.3.2 zusammengefasst sind.